\chapter{Fundamentals}
\label{ch:fundamentals}

This chapter establishes the terminology and the stable, established knowledge on
which the remainder of the thesis builds. It first describes automated driving
systems and their operating envelope, then the two safety paradigms that frame the
work, followed by safety assurance cases and the notation used to express them, the
definition of a safety performance indicator, the surrogate safety metrics that serve
as candidate signals, and the simulation environment used for evaluation. The current
state of research that builds on these foundations is reviewed in
Chapter~\ref{ch:sota}.

\section{Automated Driving Systems and the Operational Design Domain}
\label{sec:fund_ads}

An \ac{ADS} is the combination of hardware and software that performs the entire
dynamic driving task on a sustained basis. The taxonomy of SAE~J3016 distinguishes
six levels of driving automation, from no automation at Level~0 to full automation at
Level~5, and reserves the term \ac{ADS} for the systems at Levels~3 to~5 that perform
the dynamic driving task without the constant supervision of a human driver
\cite{saej3016}. The dynamic driving task is commonly decomposed into the functions of
perception, prediction, planning, and control, together with a function that monitors
the operating conditions. This functional decomposition is used throughout the thesis
because it names the tasks that any driving system must perform, independently of the
particular software architecture that performs them.

Every \ac{ADS} is specified to operate only within a defined \ac{ODD}, that is, the
set of operating conditions under which it is designed to function, including
environmental, geographical, and time-of-day constraints as well as the presence or
absence of certain traffic and road features \cite{saej3016}. The \ac{ODD} is central
to safety assurance, because a claim of safety holds only within the declared
conditions, and a departure from those conditions is itself a hazard. Thresholds on
the indicators derived in this thesis are therefore conditioned on the \ac{ODD}, so
that the value considered acceptable under one set of conditions, for example dry
weather, differs from the value considered acceptable under another, for example rain.

\section{Functional Safety and the Safety of the Intended Functionality}
\label{sec:fund_safety}

Two complementary standards frame the safety of an \ac{ADS}. ISO~26262 addresses
functional safety, defined as the absence of unreasonable risk due to hazards caused
by malfunctioning behaviour of electrical and electronic systems, and provides a risk-
based development process organised around automotive safety integrity levels
\cite{iso26262}. Its scope is the fault, that is, a component that fails to behave as
specified. ISO~21448 addresses the \ac{SOTIF}, defined as the absence of unreasonable
risk due to hazards resulting from functional insufficiencies of the intended
functionality or from reasonably foreseeable misuse, when no component has failed
\cite{iso21448}. Its scope is the situation in which every component operates as
specified yet the specification itself is inadequate, which is characteristic of the
perception and machine-learning elements of an \ac{ADS}.

The distinction matters for the present work because the two paradigms call for
different reasoning. A fault-based hazard is well captured by the failure logic of
classical reliability analysis, whereas a functional insufficiency is not, since there
is no failed component to point to. The indicators derived in this thesis are intended
to observe the developing consequences of both kinds of hazard in the behaviour of the
vehicle, which is the reason the derivation reasons about functional impossibilities
rather than component failures.

\section{Safety Assurance Cases and the Goal Structuring Notation}
\label{sec:fund_gsn}

A \ac{SAC} is a structured argument, supported by a body of evidence, that a system is
acceptably safe to operate in a defined context \cite{bishop2000methodology}. It makes
the reasoning from evidence to a top-level safety claim explicit and reviewable, rather
than leaving it implicit in a collection of documents. The \ac{GSN} is a graphical
notation for expressing such arguments and is widely used in safety-critical domains
\cite{kelly2004gsn}. Its community standard defines the core element types used in this
thesis \cite{gsncs2021}, which are summarised in Figure~\ref{fig:gsn_notation}: a goal
states a claim to be established, a strategy explains how a goal is decomposed into
sub-goals, a solution provides the evidence that supports a goal, and context,
assumption, and justification elements record the conditions under which the argument
holds. An argument is built by decomposing the top-level goal through strategies into
sub-goals until each remaining goal can be discharged by a solution.

% ---------------------------------------------------------------------------
\begin{figure}[htbp]
    \centering
    \includegraphics[width=15cm]{gsn_notation}
    \caption{Core element types of the Goal Structuring Notation used in this thesis:
    goal, strategy, solution, context, assumption, and justification, together with
    the supported-by and in-context-of connectors (based on the GSN community standard
    \cite{gsncs2021}).}
    \label{fig:gsn_notation}
\end{figure}
% ---------------------------------------------------------------------------

The \ac{GSN} is adopted here for two reasons. First, its decomposition of a top-level
goal into sub-goals produces the claim hierarchy to which safety performance indicators
are attached in Chapter~\ref{ch:methodology}. Second, a goal that is expressed as a
bare assertion can be turned into a monitorable claim by attaching an indicator to it,
which is the mechanism by which the assurance case is connected to operational data.
The assurance cases of this thesis are modelled in a model-based safety-engineering
environment that supports the construction of \ac{GSN} arguments and their linkage to
other safety artefacts.

\section{Safety Performance Indicators}
\label{sec:fund_spi}

A \ac{SPI} is a metric, supported by evidence, that uses a threshold comparison to
condition a claim within a safety assurance case; it is a metric-threshold pair that
measures a specific aspect of safety \cite{laxman2024}. Three parts therefore make up
every \ac{SPI}, and all three are needed. The first is the metric, a quantity computed
from data, such as a count, a rate, or a physical measurement. The second is the
threshold, a value that separates acceptable from unacceptable, against which the metric
is compared. The third is the claim, the specific statement in the assurance case that
the comparison is taken to support or to challenge. A metric on its own does not
constitute an \ac{SPI}, because the same numerical value can be acceptable in one
operating context and unacceptable in another; only when the metric is paired with a
threshold and tied to a claim does it carry a safety meaning
\cite{laxman2024, koopman2022howsafe}.

Metrics may be quantitative or qualitative. A quantitative metric is a measured or
counted value, for example a rate of events per distance driven, whereas a qualitative
metric records the presence or the completeness of a process step, for example whether
a required review has been carried out. Because context governs acceptability, the
threshold is frequently conditioned on the operating conditions, so that the value
regarded as acceptable in one part of the \ac{ODD} differs from that in another. The
threshold itself has two distinct roles that are developed in
Chapter~\ref{ch:methodology}: a per-event threshold defines what counts as a single
violation, and a rate-level threshold defines how often such violations may occur
before the violation rate itself is anomalous.

The property that matters most for this thesis is whether an indicator is leading or
lagging, according to whether it is observed before or after a loss event. A lagging
\ac{SPI} reports an outcome after the fact; the collision rate, the injury rate, and
the count of automation disengagements are lagging indicators of an \ac{ADS}. A leading
\ac{SPI} observes a condition that precedes a loss event and therefore offers the
opportunity to act before it occurs; the rate at which the time-to-collision falls
below a set value, the rate of activation of automatic emergency braking, and the rate
at which a following gap drops below a safe distance are leading indicators
\cite{koopman2022howsafe}. An effective programme combines both, but only leading
indicators enable a proactive response, which is why their systematic construction is
the central concern of this work.

A concrete example makes the structure explicit. Consider the assurance-case goal that
the vehicle keeps a safe separation from a leading vehicle. A candidate leading \ac{SPI}
for this goal pairs the time-to-collision metric of Section~\ref{sec:fund_metrics} with
a per-event threshold $\theta$ and a rate-level threshold $T$, and conditions the goal
in the form given in Equation~\ref{eq:spi_example}.

\begin{equation}
    \mathrm{count}\!\left(\mathrm{TTC}_{\min} < \theta\right)\ \mathrm{per}\ N\ \mathrm{manoeuvres}\ <\ T
    \label{eq:spi_example}
\end{equation}
\equationentry{Example leading safety performance indicator on time-to-collision}

where $\mathrm{TTC}_{\min}$ is the minimum time-to-collision observed during a
manoeuvre, $\theta$ is the per-event threshold below which the manoeuvre counts as a
violation, $N$ is a fixed number of manoeuvres over which the count is taken, and $T$
is the rate-level threshold. A measured violation rate at or above $T$ challenges the
safe-separation goal and, within an operational safety assurance process, triggers a
review of that part of the assurance case rather than waiting for a collision to be
recorded. The concept, its origin in the aviation safety management systems of the
International Civil Aviation Organization \cite{icao9859}, and its requirement for
automated driving by ANSI/UL~4600 \cite{ul4600} and BSI PAS~1881 \cite{pas1881} are all
introduced here as fundamentals; how a leading \ac{SPI} such as this one is derived
systematically from a claim is the subject of Chapter~\ref{ch:methodology}.

\section{Surrogate Safety Metrics}
\label{sec:fund_metrics}

Because collisions are rare, safety is commonly assessed through surrogate metrics that
quantify how close an interaction comes to a collision, and these metrics supply the
signal component of the indicators derived later. The most established is the
\ac{TTC}, introduced by Hayward, which expresses the time remaining before two road
users on a collision course would collide if their velocities remained unchanged
\cite{hayward1972ttc}. For a following vehicle approaching a leading vehicle along the
same path, the \ac{TTC} is given by Equation~\ref{eq:ttc}.

\begin{equation}
    \mathrm{TTC} = \frac{d_\mathrm{rel}}{v_\mathrm{rel}}, \qquad v_\mathrm{rel} > 0
    \label{eq:ttc}
\end{equation}
\equationentry{Time-to-collision}

where:
\begin{tabbing}
    \hspace{2.2cm} \= \kill
    $d_\mathrm{rel}$ \> relative distance between the two road users in m \\
    $v_\mathrm{rel}$ \> relative closing speed, positive when the gap is closing, in m/s
\end{tabbing}

A related time-based measure is the \ac{PET}, introduced by Allen et al., which
applies to crossing conflicts and measures the time between the moment the first road
user leaves a shared conflict area and the moment the second enters it
\cite{allen1978pet}. A small \ac{PET} indicates that two road users occupied nearly
the same space at nearly the same time. Where a required response is of interest rather
than the time available, the required deceleration to avoid a collision, given in
Equation~\ref{eq:areq}, expresses how hard the vehicle would have to brake, and a value
approaching the physical limit of the vehicle indicates a critical situation.

\begin{equation}
    a_\mathrm{req} = \frac{v_\mathrm{rel}^{2}}{2\,d_\mathrm{rel}}
    \label{eq:areq}
\end{equation}
\equationentry{Required deceleration to avoid a collision}

where:
\begin{tabbing}
    \hspace{2.2cm} \= \kill
    $a_\mathrm{req}$ \> required deceleration in m/s\textsuperscript{2} \\
    $v_\mathrm{rel}$ \> relative closing speed in m/s \\
    $d_\mathrm{rel}$ \> relative distance in m
\end{tabbing}

A formal alternative defines a safe distance from explicit assumptions about
behaviour. The Responsibility-Sensitive Safety model of Shalev-Shwartz et al. defines
the minimum safe longitudinal distance behind a leading vehicle from the response time
and the assumed acceleration and braking capabilities, as given in
Equation~\ref{eq:rss} \cite{shalev2017rss}. A distance below this value denotes a state
that the model regards as unsafe, which makes the model a natural source of a
distance-based threshold.

\begin{equation}
    d_\mathrm{min} = \left[ v_r \rho + \frac{1}{2} a_\mathrm{max}^\mathrm{acc}\,\rho^{2}
    + \frac{\left(v_r + \rho\,a_\mathrm{max}^\mathrm{acc}\right)^{2}}{2\,a_\mathrm{min}^\mathrm{brake}}
    - \frac{v_f^{2}}{2\,a_\mathrm{max}^\mathrm{brake}} \right]_{+}
    \label{eq:rss}
\end{equation}
\equationentry{Responsibility-Sensitive Safety minimum safe longitudinal distance}

where:
\begin{tabbing}
    \hspace{2.2cm} \= \kill
    $v_r,\ v_f$ \> speed of the rear and the front vehicle in m/s \\
    $\rho$ \> response time of the rear vehicle in s \\
    $a_\mathrm{max}^\mathrm{acc}$ \> maximum acceleration during the response time in m/s\textsuperscript{2} \\
    $a_\mathrm{min}^\mathrm{brake}$ \> minimum braking of the rear vehicle in m/s\textsuperscript{2} \\
    $a_\mathrm{max}^\mathrm{brake}$ \> maximum braking of the front vehicle in m/s\textsuperscript{2} \\
    $[\,\cdot\,]_{+}$ \> the positive part, that is $\max(\cdot,0)$
\end{tabbing}

The four measures differ in the conflict geometry for which each is defined and in the
quantity each returns, and Table~\ref{tab:surrogate_metrics} summarises these
differences, since the suitability of a signal for a given conflict is what governs its
selection in Chapter~\ref{ch:methodology}.

% ---------------------------------------------------------------------------
\begin{table}[htbp]
    \centering
    \caption{Surrogate safety metrics used as candidate indicator signals, with the
    conflict geometry for which each is defined, the quantity it returns, and the
    tail in which a critical value lies.}
    \label{tab:surrogate_metrics}
    \small
    \begin{tabular}{@{}p{3.5cm} p{3.3cm} p{3.4cm} p{3.2cm}@{}}
        \toprule
        \textbf{Metric} & \textbf{Conflict geometry} & \textbf{Quantity returned} & \textbf{Critical tail} \\
        \midrule
        Time-to-collision & following, or any conflict with a closing range rate
        & time remaining to contact & lower \\[2pt]
        Post-encroachment time & crossing & time between successive occupancies of a
        shared area & lower \\[2pt]
        Required deceleration & following, and turning conflicts with a defined
        stopping point & braking effort demanded & upper \\[2pt]
        Responsibility-Sensitive Safety distance & longitudinal & minimum admissible
        separation & lower \\
        \bottomrule
    \end{tabular}
\end{table}
% ---------------------------------------------------------------------------

Beyond time- and distance-based measures, model-predictive metrics such as the Model
Predictive Instantaneous Safety Metric of Weng et al. assess safety by predicting the
evolution of a scene under a control model \cite{weng2020mprism}. The geometry of the
two representative conflicts, a following conflict for time-to-collision and required
deceleration and a crossing conflict for post-encroachment time, is shown in
Figure~\ref{fig:metrics_geometry}. The suitability of each metric depends on this
geometry, and a comprehensive survey of criticality metrics and their applicability is
provided by Westhofen et al. \cite{westhofen2023}. This dependence on conflict geometry
is examined further in Chapter~\ref{ch:sota}.

% ---------------------------------------------------------------------------
\begin{figure}[htbp]
    \centering
    \includegraphics[width=15cm]{surrogate_metrics_geometry}
    \caption{Geometry of the surrogate safety metrics used in this work. Left: a
    following conflict, in which the relative distance and closing speed define the
    time-to-collision and the required deceleration. Right: a crossing conflict, in
    which the post-encroachment time is the interval between the first road user
    leaving and the second entering the shared conflict area.}
    \label{fig:metrics_geometry}
\end{figure}
% ---------------------------------------------------------------------------

\section{Simulation-Based Safety Assessment}
\label{sec:fund_sim}

Because field data on rare safety-critical events is sparse and reporting is
inconsistent, safety-critical behaviour is studied in simulation, where scenarios can
be generated in large numbers and with controlled parameters. CARLA is an open-source
simulator for automated driving research that provides sensor models, traffic
participants, and configurable environmental conditions, together with a programmable
interface for scripting scenarios and logging ground-truth state \cite{dosovitskiy2017carla}.
In this thesis it is used to generate the scenarios from which the derived indicators
are evaluated. The use of ground-truth state has an important consequence for the scope
of the evaluation, which is discussed in Chapter~\ref{ch:methodology}, because signals
that depend on the internal outputs of a perception module cannot be observed when the
ego vehicle is driven by a rule-based agent on ground truth.